\section{Background: Induced Demand and Managed Lane Design}
\label{sec:literature}

\subsection{The Fundamental Law of Road Congestion}

The most robust empirical finding in transportation economics is that new road capacity is fully utilized within a decade. Duranton and Turner \citep{Duranton2011} demonstrated this in a landmark study of the continental US road network: using instrumental variables to address endogeneity between lane construction and traffic volume, they found a long-run elasticity of vehicle kilometers traveled with respect to lane kilometers of approximately 1.0. Adding 10\% more lane-kilometers produces 10\% more vehicle-kilometers traveled, restoring the original vehicle-to-capacity ratio within the absorption period. Noland \citep{Noland2001} reached consistent findings using panel data across US metropolitan areas, with induced demand elasticities of 0.5--0.9 in the short run and approaching 1.0 over 5--10 years.

The mechanism is not mysterious: new capacity reduces travel time, which increases the attractiveness of car travel relative to alternatives; new users enter the market (latent demand), existing users shift trips to newly accessible time windows, and route diversion from parallel roads adds traffic from adjacent corridors \citep{Duranton2011}. The result is that general purpose lane additions are not permanent capacity investments --- they are temporary improvements with a defined shelf life, typically 8--12 years before the restored V/C ratio makes the investment effectively complete from a congestion standpoint.

\subsection{Why Managed Lanes Avoid Induced Demand}

Managed freight lanes escape the induced demand mechanism through access restriction. Induced demand accrues to the vehicle classes that can access the new capacity; passenger vehicles drive the dominant share of total VMT growth and therefore dominate the induced demand response. A freight-only managed lane that physically excludes passenger vehicles removes the primary induced-demand vehicle class from the new capacity pool. Induced demand on the managed lanes accrues only to freight --- which grows at 1.8\% per year per FHWA freight projections \citep{FHWA_freight2023}, compared to passenger demand growth of 2.4\% per year. At 1.8\% annual growth, the managed freight lane capacity designed to V/C $\leq$ 0.85 at opening reaches V/C = 1.0 in approximately 9 years. At 30-year design life, this implies roughly two capacity refresh cycles on the freight lanes, compared to zero on a general purpose lane addition that reaches V/C = 1.0 within a decade and provides no incremental capacity benefit thereafter.

This is not a speculative model. The I-66 managed lane program in Virginia, while not freight-exclusive, demonstrated that access-controlled lanes (requiring HOV-2 or dynamic pricing) maintained average speeds of 55+ mph during peak periods while adjacent general purpose lanes operated at LOS F \citep{FHWA_ML2021}. The access control mechanism is the operative variable: managed lanes that lose access control through enforcement failure or regulatory dilution quickly converge to general purpose lane performance.

\subsection{Managed Lane Precedents and the Novel Freight-Exclusive Contribution}

Managed lane programs in the US have historically focused on high-occupancy toll (HOT) lanes and value-priced express lanes for passenger vehicles. The SR-91 express lanes in Orange County, California, opened in 1995 as the first value-priced highway lanes in the US and demonstrated that dynamic pricing could maintain LOS B on managed lanes while adjacent free lanes operated at LOS F \citep{FHWA_ML2021}. The Houston Katy Freeway HOT lanes and the I-95 Express Lanes in South Florida followed the same model. None of these programs are freight-exclusive: they permit any vehicle willing to pay the dynamic toll, which means they are subject to induced demand from the passenger vehicle market.

The Interstate 2.0 managed freight lane proposal is the first systematic application of the freight-exclusive model to T1 corridors. The contribution is not the managed lane concept --- that is well-established --- but the application of access restriction at the vehicle-class level (freight-classified vehicles only, not price-filtered vehicles) to eliminate induced passenger demand permanently rather than managing it through pricing. The distinction matters because pricing is a flow-control mechanism, not a capacity allocation mechanism: a price-cleared managed lane in a sufficiently congested corridor will eventually find a clearing price that admits the maximum volume it can carry, restoring V/C to design capacity regardless of vehicle class composition.

\subsection{Truck Platooning and the V2X Benefit}

Truck platooning --- coordinated convoy operation of 2--4 trucks with reduced following distances --- reduces aerodynamic drag by 20--35\% for trailing vehicles in the platoon \citep{Browand2004}. The drag reduction is a function of following distance and vehicle configuration: a following truck at 4-second headway (approximately 100 meters at 65 mph) achieves approximately 10\% drag reduction; at 1-second headway (the minimum safe gap enabled by V2X cooperative adaptive cruise control), drag reduction approaches 30--35\%. At current diesel prices (\$3.40/gallon US average in 2024) and commercial vehicle mileage (6--7 mpg for loaded combination trucks), a 25\% drag reduction on the trailing vehicle reduces fuel cost by approximately \$0.08 per mile --- an annual savings of \$180 million on a T1 corridor carrying 8,000 platoon-capable freight vehicles per day \citep{ATRI_costs2024}.

Platooning is operationally feasible only on access-controlled lanes at sustained speeds of 60--70 mph. In general purpose lanes, the mix of passenger vehicles --- which lane-change and speed-differential interact with the platoon formation --- makes sustained close-following distances unsafe and legally uncertain \citep{Browand2004}. The freight-exclusive managed lane is therefore a necessary precondition for widespread commercial platooning adoption, not merely a complementary investment.

\subsection{HCM Managed Lane Capacity Standards}

Highway Capacity Manual 7th edition \citep{HCM7} assigns managed lanes a capacity premium over general purpose lanes due to the elimination of weaving and merge conflicts at access points. Basic freeway segment capacity on a managed lane is 2,400 passenger car equivalents per hour per lane (pcphpl) at design speed $\geq$ 65 mph, compared to 2,200 pcphpl for a general purpose freeway lane. The 9\% premium reflects the reduction in lane-change turbulence and merge conflict events that reduce throughput in general purpose environments. For freight-exclusive managed lanes, where the vehicle composition is 100\% combination trucks and the PCE for trucks on level terrain is 1.5, effective freight capacity per lane is 1,600 commercial vehicle equivalents per hour per lane --- still superior to the 1,267 cvphpl achievable on a general purpose lane at the same truck composition.

FHWA managed lane design guidance \citep{FHWA_ML2021} specifies access control requirements, signing standards, and enforcement infrastructure for managed lane programs. The guidance does not address freight-exclusive lanes specifically --- a gap that the I2.0 program will need to address through Federal Register rulemaking or interim guidance, as the legal framework for vehicle-class-based access restriction on interstate highways requires clarification under 23 U.S.C. 127 (vehicle weight limits) and 23 U.S.C. 301 (interstate system access standards).
